% v8 (2026-09-14): moved to Section 4 and re-aimed. It used to open the paper and argue that the
% design is closed; it now answers the question Sections 2-3 raise, which is why a per-token meter
% could not have found what selection anchoring found. The dichotomy proposition leads rather than
% concludes, and Proposition~\ref{prop:threshold} has moved to Section~\ref{sec:selection}, where
% what it certifies is the point. v7 kept as frontier_v7_2026-09-14.tex.

\section{Why no per-token budget can do this}\label{sec:frontier}

Selection anchoring buys more utility than a metered decoder for a two-thousandth of the divergence,
under a strictly stronger order. The next question is not whether to believe that but whether the
meter can be repaired to match it. It cannot, and the obstruction is one line of the chain rule.

\begin{proposition}[An affordable causal policy is the anchor almost everywhere]\label{prop:sparse}
Let $q$ be any causal policy, $q(y) = \prod_{t \le T} q_t(y_t \mid y_{<t})$, with sequence budget
$D_{\mathrm{KL}}(q \,\|\, p_s) \le K$, and write
$N_\varepsilon = \#\{t \le T : D_{\mathrm{KL}}(q_t(\cdot \mid y_{<t}) \,\|\, p_{s,t}) > \varepsilon\}$.
Then $\mathbb{E}_q[N_\varepsilon] \le K/\varepsilon$ for every $\varepsilon > 0$.
\end{proposition}

The proof is the exact chain rule
$D_{\mathrm{KL}}(q\|p_s) = \sum_t \mathbb{E}_q[D_{\mathrm{KL}}(q_t\|p_{s,t})]$ followed by Markov at
each step. Its force is the contrapositive: a policy that moves the distribution at a constant
fraction of its steps has a budget that is $\Omega(T)$, which is the dichotomy of
Section~\ref{sec:intro}. Nothing here forbids a causal policy from \emph{concentrating} a bounded budget; it forbids one from
spreading it, which is what a rate does by construction. That is the whole gap
Section~\ref{sec:selection} walks through, and it is why selection anchoring is not one option among
several: it is what is left when the per-step axis is closed.

\paragraph{The class and the rates}
Making both horns measurable takes three quantities and no attack. Write
$p_\theta \propto p_s^{1-\theta} p_r^{\theta}$ for the geometric path between the two models, so that
$\log p_\theta = \log p_s + \theta \ell - \log Z(\theta)$ with $\ell = \log p_r - \log p_s$ and
$Z(u) = \mathbb{E}_{p_s}[e^{u\ell}]$. A \emph{bucket-metered decoder} serves the largest
$\theta_t \in [0,1]$ with $c(\theta_t) \le k_t$, for $k_t$ the allowance of a token bucket of rate
$k$ started at a debt $\delta$. The charge is the R\'enyi divergence of order $\alpha \in
[1,\infty]$,
\begin{equation}
c_\alpha(\theta) \;=\; D_\alpha(p_\theta \,\|\, p_s)
 \;=\; \frac{\log Z(\alpha\theta) - \alpha \log Z(\theta)}{\alpha - 1},
\label{eq:renyi-path}
\end{equation}
whose $\alpha \to 1$ limit is $D_{\mathrm{KL}}(p_\theta \| p_s)$, the charge of the deployed
decoder, and whose $\alpha \to \infty$ limit is the pathwise charge selection anchoring satisfies
for free, so the family contains both published accounting rules and interpolates between a
guarantee about an average and one about every event. For a protected work $x$ of $N$ characters,
with $s_1,\dots,s_T$ the safe model's per-token surprisals along $x$, $N_t$ the characters consumed
after $t$ tokens and $S(x) = \sum_t s_t$, three rates govern everything below:
\begin{equation}
\underbrace{s(x) = \frac{S(x)}{N}}_{\text{surprisal rate}}, \qquad
\underbrace{k_{\mathrm{crit}}(x) = \max_{t \le T} \frac{\sum_{i \le t} s_i + \delta}{N_t}}_{\text{Loynes running maximum}}, \qquad
\underbrace{c_{\mathrm{use}} = \operatorname{med}_{y \sim p_r} \frac{\log p_r(y) - \log p_s(y)}{|y|}}_{\text{divergence rate of ordinary traffic}} .
\label{eq:rates}
\end{equation}
All three need only forward passes --- no decoding, no attack, no access to the mechanism --- so a
deployer can compute them in advance. Rates are nats per \emph{character} where safe models with
different vocabularies are compared and per \emph{token} where a rate meets a budget; both
denominators appear in Appendix~\ref{app:opening}.

The other side of the dichotomy has a price too, and it is the one that shows a certificate and a
useful mechanism are not intrinsically in conflict.
\begin{theorem}[The price of utility]\label{thm:nfl}
Let a decoder of the class above run to a per-trajectory sequence budget $K$ under any charge
$c_\alpha$, $\alpha \in [1,\infty]$, and write $q$ for the law of its output. For every bounded
utility $U$,
\begin{equation}
K \;\ge\; \Lambda^*_s\!\bigl(\mathbb{E}_q[U]\bigr), \qquad
\Lambda^*_s(u) = \sup_{\lambda}\;\bigl\{\lambda u - \log \mathbb{E}_{p_s}[e^{\lambda U}]\bigr\},
\label{eq:nfl}
\end{equation}
the Cram\'er rate function of $U$ under the safe model.
\end{theorem}

The proof is Donsker--Varadhan at the sequence level (Appendix~\ref{app:proofs}); it is a converse
holding for the optimal policy, and that optimum is not loose, the maximum reward purchasable at a
fixed KL budget having a closed form governed by a Jeffreys divergence rather than the
$\sqrt{\mathrm{KL}}$ of earlier analyses \citep{monteiropaes2026limits}. Read beside
Proposition~\ref{prop:threshold} it settles the shape of the problem: reproducing a particular work
costs $S(x)$ nats and grows with its length --- $849$ for a protected target in
Section~\ref{sec:orders} --- while buying mean utility $u$ costs $\Lambda^*_s(u)$, which for bounded
$U$ is $O(1)$ \emph{however long the sequence}. There is room for a mechanism that is both safe and
useful, and a rate cannot occupy it.

\begin{proposition}[A per-token meter pays the imitation cost, not the utility cost]\label{prop:imitation}
Let a decoder of the class run under the deployed charge $\alpha=1$ and write
$Z_T = \sum_{t \le T} D_{\mathrm{KL}}(p_{\theta_t} \| p_{s,t})$ for its realised sequence divergence.
If the bucket constraint binds on a fraction $\beta$ of steps then
$Z_T \ge (1-\beta)\sum_{t} D_{\mathrm{KL}}(p_{r,t} \| p_{s,t})$, which is linear in $T$ and does not
depend on the utility achieved. Since $\Lambda^*_s(u) \le -\log \Pr_{p_s}[U = u_{\max}]$ for every
bounded $U$, the overhead $Z_T / \Lambda^*_s(\mathbb{E}_q[U])$ is $\Omega(T)$ whenever $\beta$ is
bounded away from $1$.
\end{proposition}

This is the vacuous horn taken \emph{by construction}, not by accident: where the bucket is slack the
largest feasible tilt is $\theta_t = 1$, so the step is charged the full
$D_{\mathrm{KL}}(p_{r,t}\|p_{s,t})$ --- the cost of \emph{imitating} the risky model, which the
decoder pays whether or not imitation is what the user wanted. Both factors are \emph{measured, not
assumed}, over $16{,}500$ logged trajectories: $\beta$ falls from $0.954$ at $k=0.1$ to $0.033$ at
$k=3$ and $0.0000$ at $k=20$, while the rate it multiplies saturates at $0.857$ nats per token and
the running spend is linear in the step index at every budget (median $R^2 \ge 0.97$). Realised
spend therefore stalls at $171.30$ nats, $4.3\%$ of the allowance the certificate is written
against, while the cap quadruples --- and all four hold at three further pairs, the authors' own
$70$B among them, at rates $0.617$, $0.303$, $0.908$ (Appendix~\ref{app:proofs}).

A second, weaker boundary sits above the vacuity threshold: the rate needed to emit $x$ at all is
the Loynes running maximum of the surprisal process \citep{loynes1962stability}, so a budget
publishes a \emph{drift} where safety is a \emph{workload maximum}
(Proposition~\ref{prop:outrun}), and the two bracket a regime where the decoder protects the work and
no budget in the family can say so (Appendix~\ref{app:opening}). Both rates are measurable without
any mechanism and are separating --- over ten safe models the margin $s(x)/c_{\mathrm{use}}$ rises
monotonically with capability, correcting an earlier audit~\citep{vijayavallabh2026audit} that read
the protected rate alone (Appendix~\ref{app:scaling}). The raw material is improving; a per-token
meter cannot use it.
